\documentclass[10pt,a4paper]{article}
\usepackage[margin=2cm,top=2.5cm]{geometry}
\usepackage{parskip}
\usepackage{microtype}
\usepackage{amsmath,amssymb}
\usepackage{booktabs,array,tabularx}
\usepackage[table]{xcolor}
\usepackage{graphicx,float}
\usepackage{titlesec}
\usepackage{enumitem}
\usepackage{tcolorbox}

\titleformat{\section}{\large\bfseries}{}{0em}{}[\titlerule]
\titleformat{\subsection}{\normalsize\bfseries}{}{0em}{}
\titlespacing{\section}{0pt}{1.4ex plus .2ex}{0.5ex}
\titlespacing{\subsection}{0pt}{0.9ex}{0.3ex}

\definecolor{highgreen}{rgb}{0.14,0.46,0.18}
\definecolor{modamber}{rgb}{0.62,0.38,0.04}
\definecolor{boxblue}{rgb}{0.90,0.94,1.00}
\definecolor{boxgreen}{rgb}{0.88,0.97,0.88}
\definecolor{boxyellow}{rgb}{1.00,0.97,0.88}
\definecolor{flagred}{rgb}{0.72,0.10,0.10}

\newcommand{\seq}[1]{\texttt{#1}}
\newcommand{\hi}[1]{\textcolor{highgreen}{\textbf{#1}}}
\newcommand{\warn}[1]{\textcolor{modamber}{\textbf{#1}}}
\newcommand{\pillar}[1]{\textcolor[rgb]{0.12,0.30,0.58}{\textbf{#1}}}

\begin{document}

\begin{center}
  {\LARGE\bfseries iGEM 2026: Concrete Modelling Project Proposal}\\[4pt]
  {\large Two-Pillar Framework for the MMP-9 / HNE Wound Biosensor}\\[3pt]
  {\small\color{gray} Eton / Dry Lab}\\[1pt]
  {\small\color{gray}
    ESM-2 substrate design $\cdot$ Latent Bayesian optimisation $\cdot$
    Multi-scale biosensor physics}
\end{center}

\vspace{0.5em}\hrule\vspace{0.8em}

% -----------------------------------------------------------------------
\section{Project Overview}

The modelling project consists of two interlocking pillars that together
address both ends of the biosensor design problem.

\pillar{Pillar 1 (ML/DL)} answers: \emph{which peptide sequence should we use?}
A protein language model (ESM-2) is fine-tuned on experimental cleavage data
and used as the backbone for a Bayesian optimisation loop that generates novel,
selectivity-optimised substrate candidates outside the training distribution.

\pillar{Pillar 2 (Physics)} answers: \emph{how will that peptide perform on the
device?}
A multi-scale model of the wound biosensor translates protease concentrations in
wound fluid into a predicted electrochemical signal, via an ODE model of the wound
fluid protease network and a diffusion-reaction model of the electrode surface.

The two pillars connect at a single point: the substrate sequence output by Pillar 1
is passed into Pillar 2 as the peptide whose surface kinetics are modelled.
Together they produce a fully computational design rationale for the biosensor,
grounded in experimental MEROPS data, that wet lab then validates.

\medskip
\begin{tcolorbox}[colback=boxblue,colframe=blue!40,boxrule=0.5pt]
\textbf{Target panel (subject to wet lab confirmation): MMP-9 and HNE (Panel B)}\\[4pt]
\small
Panel B is the recommended panel (see Panel B report): both models rated HIGH,
both substrates fully supported by primary literature, no MMP-8/MMP-9
cross-reactivity issue to manage.  If wet lab moves to Panel A (three-plex),
both pillars extend trivially by adding a third protease head / ODE species.
\end{tcolorbox}

% -----------------------------------------------------------------------
\section{iGEM Judging Criteria Alignment}

The 2026 handbook lists four criteria for the Best Model award:

\begin{enumerate}[topsep=2pt,itemsep=3pt,leftmargin=*]
  \item \textbf{How impressive is the modelling?} Both pillars use
    advanced methods uncommon in iGEM: protein language model fine-tuning,
    multi-objective Bayesian optimisation, ODE systems with parameter fitting,
    and Michaelis-Menten surface kinetics.  The combination is novel.
  \item \textbf{Did it help the team understand a part, device, or system?}
    Pillar 1 determines what to synthesise; Pillar 2 predicts how it
    will behave on the electrode.  Neither is decorative.
  \item \textbf{Were real measurements used?}
    Pillar 1 trains on MEROPS experimental cleavage data.
    Pillar 2 is parameterised from published wound fluid protease concentrations
    and literature kinetic constants.  Both are validated against wet lab output.
  \item \textbf{Good example for others?}
    Full pipeline (fine-tuned weights, BO configuration, ODE parameters) is
    documented and re-runnable on any protease pair.
\end{enumerate}

The Heidelberg 2025 gold standard: model predicts optimal experimental condition,
wet lab confirms it.  Our equivalent: Pillar 1 outputs a novel substrate sequence;
Pillar 2 predicts its LOD on the biosensor; wet lab measures both cleavage
efficiency and device signal.  The model made both predictions before the
experiment ran.

% -----------------------------------------------------------------------
\section{Pillar 1: ESM-2 Fine-tuning and Latent Bayesian Optimisation}

\subsection{Stage 1a: Fine-tuning ESM-2 as a Selectivity Encoder}

ESM-2 (\seq{esm2\_t33\_650M\_UR50D}, Lin et al.\ 2022) is pre-trained on
250 million protein sequences and provides rich evolutionary context as
per-residue representations.  The 650M variant is fine-tuned on the MEROPS
MMP-9 and HNE training sets using two parallel classification heads:

\begin{equation}
  h_k = \sigma\!\left(W_k \cdot \overline{\mathrm{ESM2}}(x) + b_k\right),
  \quad k \in \{\text{MMP-9},\;\text{HNE}\}
\end{equation}

where $\overline{\mathrm{ESM2}}(x) \in \mathbb{R}^{1280}$ is the mean-pooled
final-layer representation of the 8-mer $x$.  Fine-tuning proceeds in two
phases: head-only for 10 epochs (trunk frozen), then full fine-tuning at a
low trunk learning rate ($\eta = 10^{-5}$, discriminative per layer) to
prevent catastrophic forgetting.

The fine-tuned embedding space is semantically structured by selectivity:
sequences that share protease recognition properties cluster together.
This is the latent space in which Stage 1b operates.

\textbf{Interpretability output:} integrated gradients applied to the
fine-tuned model produce per-position attribution scores
$\alpha_{k,p}$ for protease $k$ at position $p$ (P4 through P4').
If $\alpha_{\text{MMP-9},\,\text{P1}'}$ is highest among all positions,
this independently confirms the Kridel 2001 S1' pocket result.
Positions where model attribution and biochemical literature disagree
are flagged for point-mutant synthesis.

\subsection{Stage 1b: Latent Bayesian Optimisation}

A Gaussian process with Mat\'{e}rn-5/2 kernel is trained on the MEROPS
training set to map ESM-2 embeddings $z \in \mathbb{R}^{1280}$ to the
two-dimensional selectivity vector
$[s_{\text{MMP-9}},\, s_{\text{HNE}}]$.
The GP posterior $p(s \mid z)$ provides both a prediction and a
calibrated uncertainty, which the acquisition function exploits.

Multi-objective Bayesian optimisation using Expected Hypervolume
Improvement (EHVI, implemented in BoTorch) searches the fine-tuned
latent space for Pareto-optimal sequences: maximise $s_{\text{MMP-9}}$,
minimise $s_{\text{HNE}}$ (for the MMP-9 channel) and vice versa.
The search runs for 50 BO iterations over a background pool of 500k
8-mer candidates embedded by ESM-2.

\textbf{PSSM filter:} before any candidate is passed to Pillar 2 or wet
lab, it must satisfy PSSM\textsubscript{target} $>0$ and
PSSM\textsubscript{off-target} $<2$.  This applies the existing
mechanistic scorer as a biochemical gate, bridging the new generative
pipeline and the original substrate predictor.

\textbf{Positive control:} if BO rediscovers \seq{IPRTLTAA} or
\seq{AAPVSGIG} from the training set, end-to-end pipeline consistency
is confirmed before novel sequences are committed to synthesis.

\subsection{Stage 1c: Deliverables}

\begin{itemize}[topsep=2pt,itemsep=1pt,leftmargin=*]
  \item Fine-tuned ESM-2 weights and AUROC benchmark vs.\ VHSE8 + LR baseline.
  \item Per-position attribution maps per protease (comparison to literature).
  \item Pareto front in selectivity space; top 3--5 novel candidate sequences.
  \item Positive control validation (rediscovery of known substrates).
\end{itemize}

% -----------------------------------------------------------------------
\section{Pillar 2: Multi-scale Biosensor Physics Model}

\subsection{Layer 2a: Wound Fluid Protease Network (ODE)}

The wound fluid is not a simple solution of a single protease at a known
concentration.  MMP-9 exists primarily as a pro-enzyme (pro-MMP-9) and
requires activation; HNE is inhibited by $\alpha_1$-antitrypsin (A1AT)
present in wound exudate.  The effective (active, uninhibited) enzyme
concentration is what the biosensor actually measures.

A system of ODEs models the wound fluid compartment:

\begin{align}
  \frac{d[\text{pro-MMP-9}]}{dt} &= k_\text{prod} - k_\text{act}[\text{plasmin}][\text{pro-MMP-9}] - k_\text{deg}[\text{pro-MMP-9}]\\[4pt]
  \frac{d[\text{MMP-9}]}{dt}     &= k_\text{act}[\text{plasmin}][\text{pro-MMP-9}] - k_\text{inh}[\text{TIMP-1}][\text{MMP-9}] - k_\text{deg}[\text{MMP-9}]\\[4pt]
  \frac{d[\text{HNE}]}{dt}       &= k_\text{rel} - k_\text{A1AT}[\text{A1AT}][\text{HNE}] - k_\text{deg}[\text{HNE}]
\end{align}

Parameters are taken from published kinetic studies (Lijnen 2001 for
plasmin-MMP-9 activation; Werb 1997 for TIMP-1 inhibition;
Travis \& Salvesen 1983 for A1AT-HNE).  The ODE system is solved at
steady state for three wound conditions (healing, chronic non-healing,
infected) using published wound fluid protein concentration ranges
(Trengove 1999, Ladwig 2002, Watt 2009).

\textbf{Output:} the predicted range of active [MMP-9] and [HNE] that
the electrode surface actually encounters, for each wound state.
This is the realistic input to Layer 2b, replacing the naive assumption
that total and active enzyme concentrations are equal.

\subsection{Layer 2b: Electrode Surface Kinetics}

The peptide substrate from Pillar 1 is immobilised at the electrode
surface at surface density $\Gamma$ (mol/cm$^2$).  When protease
arrives by diffusion, cleavage follows Michaelis-Menten kinetics
at the surface:

\begin{equation}
  \frac{d\Gamma_\text{cleaved}}{dt}
  = \frac{k_{cat}\,[E]\,\Gamma}
         {K_m + \Gamma / \delta}
\end{equation}

where $[E]$ is the free enzyme concentration from Layer 2a, $\delta$
is the effective diffusion layer thickness (estimated from Cottrell
equation for the electrode geometry), and $k_{cat}$, $K_m$ are
taken from literature for the chosen substrate.

A 1D diffusion-reaction PDE for the enzyme concentration profile
near the electrode surface is solved numerically (finite difference,
50-node spatial grid):

\begin{equation}
  \frac{\partial [E]}{\partial t}
  = D_E \frac{\partial^2 [E]}{\partial x^2}
  - \frac{k_{cat}\,[E]\,\Gamma}{K_m + [E]}
\end{equation}

\textbf{Sensitivity analysis:} $\Gamma$ and the linker length (which
affects the pre-exponential electron transfer factor) are swept over
physically realistic ranges to identify the surface density and linker
length that maximise signal-to-noise at the clinically relevant
protease concentrations identified in Layer 2a.  This directly
informs electrode fabrication.

\textbf{Predicted LOD:} the model outputs a predicted limit of detection
for each wound state, expressed in nM of active enzyme.  This is a
concrete, falsifiable prediction for wet lab to test.

\subsection{Layer 2c: Deliverables}

\begin{itemize}[topsep=2pt,itemsep=1pt,leftmargin=*]
  \item Steady-state active enzyme concentrations for three wound states
    (healing, chronic, infected) with uncertainty ranges from parameter variation.
  \item Optimal electrode surface density $\Gamma^*$ and linker length recommendation.
  \item Predicted LOD (nM active enzyme) for MMP-9 and HNE channels.
  \item Sensitivity analysis figures showing how LOD depends on $\Gamma$ and $k_{cat}$.
\end{itemize}

% -----------------------------------------------------------------------
\section{How the Two Pillars Connect}

\begin{center}
\small
\begin{tabular}{ll}
  \toprule
  Pillar 1 output & Pillar 2 input\\
  \midrule
  Novel substrate sequence (e.g.\ \seq{IPRTLTAA} or BO-generated) &
    $k_{cat}$, $K_m$ for chosen sequence used in Layer 2b\\
  Selectivity score (MMP-9 vs.\ HNE) &
    Confirms the sequence is not a confound in the LOD model\\
  PSSM score &
    Biochemical plausibility check before LOD prediction\\
  \bottomrule
\end{tabular}
\end{center}

\medskip
The narrative is linear and checkable: Pillar 1 picks the peptide;
Pillar 2 predicts what signal that peptide will produce on the device
at wound fluid protease concentrations.  Wet lab then measures both
the cleavage efficiency of the peptide and the device signal.
Two predictions, one experiment.

% -----------------------------------------------------------------------
\section{Open Question: Wet Lab Validation Form}

\warn{This section requires confirmation from the wet lab team before
the modelling plan is finalised.}

The specific form of wet lab validation determines which model
predictions are directly testable within the 3-week experimental
window.  Two scenarios:

\begin{enumerate}[topsep=2pt,itemsep=4pt,leftmargin=*]
  \item \textbf{FRET cleavage assay on peptides}: measures
    $k_{cat}/K_m$ for each candidate sequence.  Validates Pillar 1
    directly (is the BO-generated sequence actually cleaved?).
    Can be done on synthesised FRET peptides without a full device.
  \item \textbf{Electrochemical measurements on functionalised
    electrode}: measures device signal at known protease concentrations.
    Validates Pillar 2 directly (does the LOD prediction match?).
    Requires a working electrode and immobilised peptide.
\end{enumerate}

Ideally both are done; given 3 weeks, the FRET assay (scenario 1)
is lower barrier and directly validates the highest-value output of
the project (the novel substrate candidate).  The recommendation is to
prioritise FRET validation of the top BO-generated sequence and one
known positive control (\seq{IPRTLTAA}).

% -----------------------------------------------------------------------
\section{Summary}

\begin{center}
\small
\begin{tabular}{llll}
  \toprule
  Component & Methods & Key output & Judging criterion\\
  \midrule
  ESM-2 fine-tuning      & Transfer learning, multi-task & Selectivity encoder & 01, 03\\
  Attention analysis     & Integrated gradients          & Position importance maps & 02, 04\\
  Latent BO              & GP + EHVI, BoTorch            & Novel substrate candidates & 01, 02\\
  PSSM filter            & Log-odds (existing)           & Biochemical gate & 02, 03\\
  Wound fluid ODE        & Nonlinear ODE system          & Active enzyme concentrations & 01, 02, 03\\
  Electrode kinetics PDE & 1D diffusion-reaction         & Predicted LOD, optimal $\Gamma$ & 01, 02\\
  Wet lab validation     & FRET / electrochemical (TBD)  & Experimental confirmation & 03\\
  \bottomrule
\end{tabular}
\end{center}

\medskip
The project satisfies the Heidelberg 2025 gold standard: two concrete,
falsifiable predictions (novel substrate from Pillar 1; LOD from Pillar 2)
that wet lab tests before the wiki freeze.  The two-pillar structure
gives judges a complete model story from molecular sequence design
through to clinical signal output.

\end{document}
